\scanpage{41}
% L1–6
$A$ anneau local, quotient d'un régulier $A'$. $n=\dim.\,\mathrm{Krull}\ A$.

$\underline{k}=A/\unsure{\mathfrak r|\mathfrak m}(A)$

$\unsure{\mathfrak n|\mathfrak N}$ un idéal de définition.

$(x_1,\ldots,x_n)$ un système de paramètres.

$\mathfrak p$ un idéal premier de dim $n$.

\uline{Conditions}

% L7–9
(i) $\operatorname{Ext}^{n}(\underline k,M)\simeq k,\ \deleted{[}\operatorname{Ext}^{n+1}(\underline k,M)=\deleted{[}\operatorname{Ext}^{n-1}(\underline k,M)=0\deleted{]]}$
% slanted
\margin{\unsure{on peut coupler|on peut compléter} comme on veut !}
\margin{$\operatorname{Ext}^{n}$ et $\operatorname{Ext}^{n+1}$ ne suffisent pas ni $\operatorname{Ext}^{n}$ et $\operatorname{Ext}^{n-1}$ ni $\operatorname{Ext}^{n-1}$ et $\operatorname{Ext}^{\unsure{n+1|n-1}}$. Contre-exples $A/\mathfrak p$, $A$ et $A/\mathfrak p$ rég. $A/J$, $\Omega_{A}$, où $A$ de dim $1$ non équ. $A^m$, où $A$ rég.}

(ii) $\operatorname{Ext}^{n}(\underline k,M)\simeq k$ et $M$ \uline{fidèle} [\deleted{\unsure{ou}} ou \unsure{seult|seulem$^{\mathrm t}$} de support tt]

% L10–13
\deleted{(iii) $\operatorname{Ext}^{n+1}(\underline k,M)=\operatorname{Ext}^{n-1}(\underline k,M)=0,\ 0<\lg_{\mathfrak p}(M)\leq\lg_{\mathfrak p}A$}

\deleted{(iii bis) $\operatorname{Ext}^{\unsure{n-1|n}}(\underline k,M)=0$ $\left.\begin{smallmatrix}M\text{ fidèle}\\ A\text{ équidim}\end{smallmatrix}\right\}$ $0<\lg_{\mathfrak p}(M)\leq\lg_{\mathfrak p}(A)$}

\deleted{(iv) \unsure{bis} $\operatorname{Ext}^{n-1}(\underline k,M)=0$ et $\operatorname{Ext}^{n}(A/\mathfrak n,M)\simeq A/\mathfrak n$}

% L14–15
(ii bis) $\operatorname{Ext}^{n}(A/\mathfrak n,M)\simeq\widehat{A/\mathfrak n}$ et $M$ \uline{fidèle} [faut-il $\operatorname{Ext}^{n-1}(k,M)=0$ ? et $A$ équidim ?]

\inserted{suffit il \unsure{même|$=$} que $\operatorname{Ext}^{n}(A/\mathfrak n,M)$ soit \uline{comonogène} ?}

% L16–20
(iii)
\diagram{\scriptsize\begin{tikzcd}[ampersand replacement=\&,column sep=small,row sep=.25em]
 {} \& {\deleted{M\text{ fidèle},\, A\text{ équidim}}}
 \arrow[d,bend left=40]\arrow[ddd,bend left=48] \& {\deleted{\text{\textcircled{a}}}\ \deleted{\unsure{\text{b}}}}\\
 {\left\{\begin{array}{l}\operatorname{Ext}^{n-1}(\underline k,M)=0,\\ \text{et l'une des cond. équivalentes}\end{array}\right.}
 \& \operatorname{Ext}^{n+1}(\underline k,M)=0 \& \text{\textcircled{a}}\\
 |[draw,rounded corners=1em]|{0<\lg_{\mathfrak p}(M)\leq\lg_{\mathfrak p}(A)} \arrow[ur,no head]
 \& H^n(M)\text{ dualisant} \& \text{\textcircled{b}}\\
 {} \& H^n(M)\text{ injectif} \& \text{\textcircled{c}}
\end{tikzcd}}
% slanted
\margin{$H^n(M)$ \unsure{tt seul|tout seul} ne donne rien, car on peut rajouter à $M$ des modules de dim. $<n$, sans changer $H^n(M)$. De $=$, $H^n(M)$ et $M$ équidim \unsure{ne nous disent rien de nouveau}, car on peut \unsure{trouver des exemples}, \unsure{prendre} $M'\subset M$ où $\dim M/M'<n-1$. Alors $\operatorname{Ext}^{n-1}(k,M)=0$ est pratiquement \unsure{$\Omega_{A/J}^{m}$|comme $\Omega_{A/J}^{m}$|$\Omega_{A/\mathfrak I}$} \illegible{2} $M=A/J$ \unsure{pour $A/J$} \illegible{3}.}

% L21–25
(iv) Le foncteur $N\mapsto\operatorname{Ext}^{n}(N,M)$ est dualisant

(iv bis) Le foncteur $N\mapsto\operatorname{Ext}^{n}(N,M)$ est exact, $0<\lg_{\mathfrak p}(M)\leq\lg_{\mathfrak p}(A)$

\hfill ($N$ de long. finie)

$\uparrow$

(v) L'annulateur de $\mathfrak m$ dans $M/\sum x_iM$ est de dim $1$, et $M$ est \uline{fidèle}

$\boxed{\text{(vi)}}$ $A$ est C-M., et $M$ est $\simeq\Omega_A$

% L26–29
\margin{\uline{Pb} $\left.\begin{array}{l} M\ne0\\ \operatorname{Ext}^{n-1}(\underline k,M)=\left\{\begin{array}{l}\operatorname{Ext}^{n+1}(\underline k,M)=0\\ \text{ou }H^n(M)\text{ injectif}\end{array}\right.\end{array}\right.$ $\Longrightarrow A$ est C.M. et $M\simeq\Omega_A^m$.}
% slanted
\margin{$M\ne0$, $\operatorname{Ext}^{n}(N,M)$ exact $\Longrightarrow$}
% slanted
\margin{\unsure{trouver une condition|trouvons une condition|mais une condition} \unsure{multiplicative|multiple} \illegible{2} \unsure{caractérisant|effectivement} \unsure{les multiples} de $\Omega_{\unsure{A}}$ \unsure{$A^m$}.}

% L30–36
\deleted{Le seul \unsure{cas connu|cas} où $\operatorname{Ext}^{n-1}(k,M)\ne0$, \inserted{$M\ne0$,} est celui où $A/J$ (l'anneau équidim. associé à $A$) est C.M. et $M\simeq\Omega_A^m\times\unsure{A/\mathfrak n^{m'}|A/J^{m'}}$ [cela implique $\operatorname{Ext}^{n}(k,M)=k$ \unsure{si|s.} $n=1$].}

\deleted{\deleted{\unsure{Le seul cas}} Le seul cas connu où $\operatorname{Ext}^{n}(k,M)\simeq k$ est celui où $A/J$ est Coh-Mac et où $M=\Omega_{A/\mathfrak n}$, où $\mathfrak n$ est un idéal de $A$ tel que $A/\mathfrak n$ soit C-M. / ou, pour $n=0$, un sous-module de $\widehat A$.}

\margin{cela implique-t-il $M$ \deleted{\unsure{quasi équidim|équidim}} ?}

% archive sheet number: 40
